\documentclass[aspectratio=169]{beamer}

% Tema UFS --- Universidade Federal de Sergipe
\usetheme{UFS}

% Pacotes base
\usepackage[utf8]{inputenc}
\usepackage[portuguese]{babel}
\usepackage{amsmath,amssymb}
\usepackage{graphicx}
\usepackage{booktabs}
\usepackage{xcolor}

% TikZ e bibliotecas
\usepackage{tikz}
\usetikzlibrary{shapes.geometric, arrows.meta, positioning, matrix, fit,
  backgrounds, decorations.pathreplacing, shadows, patterns, calc}

% Listagens --- paleta VS Code Light+
\usepackage{listings}
\definecolor{bgcolor}{HTML}{FFFFFF}
\definecolor{linecolor}{HTML}{DDDDDD}
\definecolor{numcolor}{HTML}{999999}
\definecolor{kwcolor}{HTML}{0000FF}
\definecolor{strcolor}{HTML}{A31515}
\definecolor{cmtcolor}{HTML}{008000}

\lstdefinestyle{ufsStyle}{
  backgroundcolor=\color{bgcolor},
  basicstyle=\ttfamily\footnotesize\color{black},
  keywordstyle=\color{kwcolor}\bfseries,
  stringstyle=\color{strcolor},
  commentstyle=\color{cmtcolor}\itshape,
  numberstyle=\tiny\color{numcolor},
  numbers=left, numbersep=12pt,
  xleftmargin=20pt, framexleftmargin=20pt,
  breaklines=true, tabsize=4,
  keepspaces=true, showspaces=false,
  showstringspaces=false, showtabs=false,
  frame=single, rulecolor=\color{linecolor},
  framesep=4pt, captionpos=b, upquote=true,
  columns=fullflexible,
}
\lstset{style=ufsStyle, language=C}

% --- Estilos TikZ reutilizados ---
\tikzset{
  cel/.style={rectangle, draw=UFSBlue, fill=UFSBlue!8, minimum width=1.1cm,
    minimum height=0.75cm, font=\small\ttfamily},
  celp/.style={rectangle, draw=UFSRed, fill=UFSRed!8, minimum width=1.1cm,
    minimum height=0.75cm, font=\small\ttfamily},
  celv/.style={rectangle, draw=UFSGray, fill=black!5, minimum width=1.1cm,
    minimum height=0.75cm, font=\small\ttfamily},
  rot/.style={font=\scriptsize\ttfamily, text=UFSGray},
  leg/.style={font=\scriptsize, text=UFSGray},
  seta/.style={->, thick, UFSBlue},
  setar/.style={->, thick, UFSRed},
  caixa/.style={rectangle, draw=UFSGray, rounded corners, fill=black!3,
    font=\scriptsize, align=center, inner sep=4pt},
}

% --- Metadados ---
\title{Alocação dinâmica de memória}
\subtitle{Quando o tamanho só se sabe rodando: \texttt{malloc}, \texttt{calloc}, \texttt{realloc} e \texttt{free}}
\author{Prof.\ André Yoshiaki Kashiwabara}
\institute{DCOMP -- Universidade Federal de Sergipe\\
\texttt{andreyoshiaki@dcomp.ufs.br}}
\date{}

\begin{document}

\begin{frame}[plain]
  \titlepage
\end{frame}

\begin{frame}{O que você vai aprender hoje}
  \begin{center}
  \begin{tikzpicture}[
    item/.style={rectangle, draw=UFSBlue, fill=UFSBlue!10, rounded corners,
      minimum width=10cm, minimum height=0.75cm, font=\small}
  ]
    \node[item] (t1) at (0,0)     {Por que \texttt{MAX\_CONTATOS} é uma decisão que você não deveria tomar};
    \node[item] (t2) at (0,-0.95) {\texttt{malloc} e \texttt{free}: pedir um bloco e devolvê-lo};
    \node[item] (t3) at (0,-1.90) {\texttt{calloc}: o mesmo pedido, com garantia de zeros};
    \node[item] (t4) at (0,-2.85) {\texttt{realloc}: mudar de ideia sobre o tamanho};
    \node[item] (t5) at (0,-3.80) {As quatro falhas clássicas --- e como o compilador as denuncia};
    \draw[seta] (t1)--(t2);
    \draw[seta] (t2)--(t3);
    \draw[seta] (t3)--(t4);
    \draw[seta] (t4)--(t5);
  \end{tikzpicture}
  \end{center}
  \begin{center}
  \small Hoje o ponteiro deixa de apontar para o vetor de outra pessoa:
  ele aponta para memória que \emph{você} pediu.
  \end{center}
\end{frame}

% ==================================================================
\section{Por que alocar em tempo de execução}
% ==================================================================

\begin{frame}{O problema com \texttt{MAX\_CONTATOS}}
  \begin{columns}[T]
    \begin{column}{0.5\textwidth}
      \textbf{O que fizemos até agora}\\[0.3em]
      \texttt{char nomes[4][20];}

      \medskip
      O tamanho é decidido por você, ao escrever o programa. O compilador
      reserva o espaço antes de o programa rodar.
    \end{column}
    \begin{column}{0.5\textwidth}
      \textbf{Os dois modos de errar}\\[0.3em]
      \textcolor{UFSRed}{Pequeno demais:} chega o 5º contato e não há para
      onde ir.

      \medskip
      \textcolor{UFSRed}{Grande demais:} \texttt{nomes[10000][20]} reserva
      200~KB para guardar três nomes.

      \medskip
      Quem sabe o tamanho certo é o \emph{usuário}, enquanto o programa roda.
    \end{column}
  \end{columns}
\end{frame}

\begin{frame}{Onde a memória do seu programa vive}
  \begin{center}
  \begin{tikzpicture}[node distance=0pt]
    \node[caixa, minimum width=4.4cm, minimum height=1.5cm, fill=UFSBlue!8,
          draw=UFSBlue] (pilha)
      {\textbf{Pilha (automática)}\\[2pt]
       \texttt{int x;} \quad \texttt{char nome[20];}\\[2pt]
       nasce e morre com a função\\
       tamanho fixo, decidido na compilação};
    \node[caixa, minimum width=4.4cm, minimum height=1.5cm, fill=UFSRed!8,
          draw=UFSRed, right=1.6cm of pilha] (heap)
      {\textbf{Heap (dinâmica)}\\[2pt]
       \texttt{malloc}, \texttt{calloc}, \texttt{realloc}\\[2pt]
       vive até \emph{você} chamar \texttt{free}\\
       tamanho decidido rodando};
    \draw[seta] (pilha) -- node[above, leg] {um ponteiro} node[below, leg] {liga os dois} (heap);
  \end{tikzpicture}
  \end{center}
  \begin{block}{A regra que muda tudo}
    Na pilha, quem cria também destrói: a função termina, o espaço some.
    Na heap, o espaço permanece até você devolvê-lo --- e permanece
    \emph{mesmo} que você perca o ponteiro para ele.
  \end{block}
\end{frame}

\begin{frame}[fragile]{O programa de hoje: \texttt{agenda.c} (1/4)}
\begin{lstlisting}[basicstyle=\ttfamily\scriptsize]
#include <stdio.h>
#include <stdlib.h>   /* malloc, calloc, realloc, free */
#include <string.h>

int main(void) {
    /* o que vem dentro de um bloco recem-alocado? */
    int *sujo = malloc(4 * sizeof(int));
    for (int i = 0; i < 4; i++)
        sujo[i] = 777;
    free(sujo);
    int *a = malloc(4 * sizeof(int));
    int *b = calloc(4, sizeof(int));
    printf("A: malloc -> %d %d %d %d\n", a[0], a[1], a[2], a[3]);
    printf("B: calloc -> %d %d %d %d\n", b[0], b[1], b[2], b[3]);
    free(a);
    free(b);
\end{lstlisting}
\end{frame}

\begin{frame}[fragile]{\texttt{agenda.c} (2/4): uma agenda sem \texttt{MAX\_CONTATOS}}
\begin{lstlisting}[basicstyle=\ttfamily\scriptsize, firstnumber=18]
    int capacidade = 2, n = 0;
    char **nomes  = malloc(capacidade * sizeof(char *));
    int   *idades = malloc(capacidade * sizeof(int));
    if (nomes == NULL || idades == NULL) {
        printf("sem memoria\n");
        return 1;
    }

    const char *entrada[] = {"ana", "bruno", "carla",
                             "davi", "elis"};
    int idade[] = {23, 19, 31, 24, 28};
\end{lstlisting}
\end{frame}

\begin{frame}[fragile]{\texttt{agenda.c} (3/4): o que fazer quando enche}
\begin{lstlisting}[basicstyle=\ttfamily\scriptsize, firstnumber=30]
    for (int i = 0; i < 5; i++) {
        if (n == capacidade) {
            int nova = capacidade * 2;
            char **n2 = realloc(nomes,  nova * sizeof(char *));
            int   *i2 = realloc(idades, nova * sizeof(int));
            if (n2 == NULL || i2 == NULL) { return 1; }
            printf("C: cheia com %d -> cresce para %d | nomes: %s\n",
                   capacidade, nova,
                   (n2 == nomes) ? "mesmo bloco" : "bloco NOVO");
            nomes = n2;  idades = i2;  capacidade = nova;
        }
        nomes[n] = malloc(strlen(entrada[i]) + 1);
        strcpy(nomes[n], entrada[i]);
        idades[n] = idade[i];
        n++;
    }
\end{lstlisting}
\end{frame}

\begin{frame}[fragile]{\texttt{agenda.c} (4/4): mostrar e devolver}
\begin{lstlisting}[basicstyle=\ttfamily\scriptsize, firstnumber=47]
    printf("D: %d contatos guardados em capacidade %d\n",
           n, capacidade);
    for (int i = 0; i < n; i++)
        printf("   %-6s %3d\n", nomes[i], idades[i]);
    printf("E: sizeof(nomes) = %zu | strlen(nomes[1]) + 1 = %zu\n",
           sizeof(nomes), strlen(nomes[1]) + 1);

    for (int i = 0; i < n; i++)
        free(nomes[i]);
    free(nomes);
    free(idades);
    return 0;
}
\end{lstlisting}
\end{frame}

\begin{frame}{O que este programa imprime?}
  \begin{block}{Antes de compilar, escreva sua resposta}
    \begin{enumerate}
      \item Na linha \textbf{A}, o que sai de um bloco de \texttt{malloc} que
        ninguém preencheu ainda?
      \item Na linha \textbf{B}, o que muda por ter usado \texttt{calloc}?
      \item Quantas vezes a linha \textbf{C} aparece, se a capacidade começa em
        2 e chegam 5 contatos?
      \item Na linha \textbf{C}, o vetor continua no \emph{mesmo bloco} ou vai
        para um \emph{bloco novo}?
      \item Na linha \textbf{E}, quanto vale \texttt{sizeof(nomes)}, sendo
        \texttt{nomes} uma agenda de 8 posições?
    \end{enumerate}
  \end{block}
\end{frame}

\begin{frame}[fragile]{Compile, execute e compare}
\begin{lstlisting}[basicstyle=\ttfamily\scriptsize, numbers=none, language=bash]
$ gcc -Wall -Wextra agenda.c -o agenda
$ ./agenda
A: malloc -> 0 0 0 0
B: calloc -> 0 0 0 0
C: cheia com 2 -> cresce para 4 | nomes: bloco NOVO
C: cheia com 4 -> cresce para 8 | nomes: bloco NOVO
D: 5 contatos guardados em capacidade 8
   ana     23
   bruno   19
   carla   31
   davi    24
   elis    28
E: sizeof(nomes) = 8 | strlen(nomes[1]) + 1 = 6
\end{lstlisting}
\end{frame}

% ==================================================================
\section{\texttt{malloc} e \texttt{free}}
% ==================================================================

\begin{frame}{Definição: \texttt{malloc}}
  \begin{block}{\texttt{void *malloc(size\_t n);}}
    Pede ao sistema um bloco de \textbf{n bytes} contíguos na heap e devolve o
    endereço do primeiro byte --- ou \texttt{NULL} se não houver memória.
  \end{block}
  \begin{exampleblock}{Intuição}
    \texttt{malloc} não sabe o que você vai guardar: ele conta \emph{bytes},
    não elementos. Por isso o pedido quase sempre tem a forma
    \texttt{quantidade * sizeof(tipo)}.

    \medskip
    \texttt{int *v = malloc(100 * sizeof(int));} \quad --- 100 inteiros

    \texttt{int *v = malloc(100);} \quad --- 100 \emph{bytes}: 25 inteiros
  \end{exampleblock}
\end{frame}

\begin{frame}{Anatomia de uma alocação}
  \begin{center}
  \begin{tikzpicture}
    \node[cel, minimum width=2.2cm] (p) at (0,0) {0x16f\ldots};
    \node[leg, above=2pt of p] {\texttt{int *v} --- na pilha};
    \node[leg, below=2pt of p] {8 bytes: só um endereço};

    \node[cel] (c0) at (4.2,0) {?};
    \node[cel, right=0pt of c0] (c1) {?};
    \node[cel, right=0pt of c1] (c2) {?};
    \node[cel, right=0pt of c2] (c3) {?};
    \node[leg, above=2pt of c1.north east] {bloco de \texttt{4 * sizeof(int)} = 16 bytes --- na heap};
    \node[rot, below=2pt of c0] {v[0]};
    \node[rot, below=2pt of c1] {v[1]};
    \node[rot, below=2pt of c2] {v[2]};
    \node[rot, below=2pt of c3] {v[3]};
    \draw[seta] (p) -- (c0);
  \end{tikzpicture}
  \end{center}
  \begin{block}{Duas coisas separadas}
    O ponteiro é uma variável comum, na pilha, que morre com a função.
    O bloco é outra coisa, na heap, que \textbf{não} morre com a função.
    Perder o ponteiro não apaga o bloco --- só torna impossível devolvê-lo.
  \end{block}
\end{frame}

\begin{frame}[fragile]{O contrato: sempre verifique o \texttt{NULL}}
\begin{lstlisting}[basicstyle=\ttfamily\scriptsize]
int *v = malloc(n * sizeof(int));
if (v == NULL) {              /* o pedido pode falhar */
    printf("sem memoria\n");
    return 1;
}
/* daqui para baixo, v e' um vetor de n inteiros */
\end{lstlisting}
  \begin{alertblock}{Por que isso importa}
    Se \texttt{malloc} falha e você não testa, \texttt{v[0] = 3} escreve no
    endereço 0: o programa morre com \emph{segmentation fault} em um lugar que
    nada tem a ver com a causa.
  \end{alertblock}
\end{frame}

\begin{frame}{O ciclo de vida de um bloco}
  \begin{center}
  \begin{tikzpicture}[node distance=0.9cm]
    \node[caixa, fill=UFSBlue!10, draw=UFSBlue, minimum width=2.6cm] (a) {\textbf{alocar}\\ \texttt{malloc} / \texttt{calloc}};
    \node[caixa, right=of a, minimum width=2.6cm] (t) {\textbf{testar}\\ \texttt{if (p == NULL)}};
    \node[caixa, right=of t, minimum width=2.6cm] (u) {\textbf{usar}\\ \texttt{p[i] = \ldots}};
    \node[caixa, right=of u, fill=UFSRed!10, draw=UFSRed, minimum width=2.6cm] (f) {\textbf{liberar}\\ \texttt{free(p)}};
    \draw[seta] (a)--(t);  \draw[seta] (t)--(u);  \draw[seta] (u)--(f);
    \draw[setar, dashed] (f.south) .. controls +(0,-0.8) and +(0,-0.8) .. (u.south)
      node[midway, below, leg] {depois do \texttt{free}, \texttt{p} não serve mais};
  \end{tikzpicture}
  \end{center}
  \begin{block}{Um \texttt{free} para cada \texttt{malloc}}
    Cada chamada bem-sucedida de \texttt{malloc}, \texttt{calloc} ou
    \texttt{realloc} corresponde a exatamente uma chamada de \texttt{free}.
    Nem zero, nem duas.
  \end{block}
\end{frame}

\begin{frame}{O que \texttt{free} faz --- e o que não faz}
  \begin{columns}[T]
    \begin{column}{0.5\textwidth}
      \textbf{Faz}
      \begin{itemize}
        \item devolve o bloco ao sistema, que pode reaproveitá-lo no próximo
          \texttt{malloc};
        \item aceita \texttt{free(NULL)} sem reclamar.
      \end{itemize}
    \end{column}
    \begin{column}{0.5\textwidth}
      \textbf{Não faz}
      \begin{itemize}
        \item não altera o ponteiro: \texttt{p} continua com o mesmo endereço,
          agora inválido;
        \item não apaga o conteúdo: os bytes antigos podem continuar lá;
        \item não avisa você se esquecer de chamá-lo.
      \end{itemize}
    \end{column}
  \end{columns}
  \begin{exampleblock}{Hábito barato}
    \texttt{free(p); p = NULL;} --- um ponteiro nulo falha alto e cedo;
    um ponteiro pendurado falha baixo e tarde.
  \end{exampleblock}
\end{frame}

\begin{frame}{Recapitulando: \texttt{malloc} e \texttt{free}}
  \begin{block}{O que aprendemos}
    \begin{itemize}
      \item \texttt{malloc(n * sizeof(tipo))} devolve o endereço de um bloco de
        bytes na heap, ou \texttt{NULL};
      \item o ponteiro vive na pilha; o bloco vive na heap e sobrevive à função;
      \item testar o \texttt{NULL} é parte da alocação, não um extra;
      \item \texttt{free} devolve o bloco, mas não mexe no seu ponteiro.
    \end{itemize}
  \end{block}
\end{frame}

% ==================================================================
\section{\texttt{calloc}}
% ==================================================================

\begin{frame}{Por que a linha A imprimiu \texttt{0 0 0 0}?}
  \begin{block}{O que o programa fez}
    Preencheu um bloco com \texttt{777}, liberou o bloco e pediu outro do mesmo
    tamanho --- que, muito provavelmente, é o mesmo bloco de volta.
  \end{block}
  \begin{columns}[T]
    \begin{column}{0.52\textwidth}
      \textbf{A resposta depende da máquina:}
      \begin{itemize}
        \item nesta aqui: \texttt{0 0 0 0}
        \item em muitos Linux: \texttt{777 777 777 777}
        \item com \texttt{MallocScribble=1}: \texttt{-1431655766}
      \end{itemize}
    \end{column}
    \begin{column}{0.48\textwidth}
      \textbf{Três respostas para o mesmo código.}

      \medskip
      Isso não é um detalhe de implementação curioso: é a definição de
      \emph{valor indeterminado}.
    \end{column}
  \end{columns}
\end{frame}

\begin{frame}{A regra: zero por sorte não é zero por garantia}
  \begin{alertblock}{O conteúdo de um bloco de \texttt{malloc} é indeterminado}
    Ler \texttt{v[i]} antes de escrever nele é comportamento indefinido. Se hoje
    saiu zero, foi coincidência: o sistema entregou uma página limpa ou um bloco
    que por acaso estava zerado. Amanhã, com outra entrada, sai outra coisa.
  \end{alertblock}
  \begin{exampleblock}{Consequência prática}
    Um contador que começa em \texttt{v[i]++} sem inicialização passa nos seus
    testes e falha na máquina do professor. É exatamente o tipo de bug que só
    aparece longe de quem o escreveu.
  \end{exampleblock}
\end{frame}

\begin{frame}{Definição: \texttt{calloc}}
  \begin{block}{\texttt{void *calloc(size\_t quantidade, size\_t tamanho);}}
    Aloca espaço para \texttt{quantidade} elementos de \texttt{tamanho} bytes
    cada e \textbf{preenche tudo com zeros} antes de devolver o endereço.
  \end{block}
  \begin{center}
  \small
  \begin{tabular}{lll}
    \toprule
     & \texttt{malloc} & \texttt{calloc} \\
    \midrule
    argumentos      & 1: total de bytes            & 2: quantidade e tamanho \\
    conteúdo        & indeterminado                & todos os bytes em zero \\
    custo           & só reserva                   & reserva e limpa \\
    \texttt{n * sz} transborda & silenciosamente   & detectado, devolve \texttt{NULL} \\
    \bottomrule
  \end{tabular}
  \end{center}
\end{frame}

\begin{frame}[fragile]{Quando usar cada um}
\begin{lstlisting}[basicstyle=\ttfamily\scriptsize]
/* vou preencher tudo agora mesmo: zerar seria trabalho jogado fora */
int *copia = malloc(n * sizeof(int));
for (int i = 0; i < n; i++) copia[i] = origem[i];

/* vou acumular: preciso comecar do zero */
int *contagem = calloc(256, sizeof(int));
for (int i = 0; texto[i] != '\0'; i++) contagem[(unsigned char) texto[i]]++;
\end{lstlisting}
  \begin{exampleblock}{Regra de bolso}
    Se a primeira coisa que você faria é um laço escrevendo zeros, use
    \texttt{calloc}. Se a primeira coisa é um laço escrevendo os valores de
    verdade, use \texttt{malloc}.
  \end{exampleblock}
\end{frame}

\begin{frame}{Recapitulando: \texttt{calloc}}
  \begin{block}{O que aprendemos}
    \begin{itemize}
      \item \texttt{malloc} entrega bytes com conteúdo indeterminado --- zeros
        observados são coincidência, não contrato;
      \item \texttt{calloc(q, t)} entrega \texttt{q * t} bytes zerados e ainda
        protege contra o transbordo da multiplicação;
      \item zerar custa tempo: escolha pelo que você vai fazer com o bloco, não
        por hábito.
    \end{itemize}
  \end{block}
\end{frame}

% ==================================================================
\section{\texttt{realloc}}
% ==================================================================

\begin{frame}{Por que a agenda mudou de bloco duas vezes?}
  \begin{block}{A saída}
    \texttt{C: cheia com 2 -> cresce para 4 | nomes: bloco NOVO}\\
    \texttt{C: cheia com 4 -> cresce para 8 | nomes: bloco NOVO}
  \end{block}
  \begin{columns}[T]
    \begin{column}{0.5\textwidth}
      \textbf{A pergunta}\\[0.3em]
      Se o bloco é novo, o que aconteceu com os nomes que já estavam lá?

      E com o bloco antigo?
    \end{column}
    \begin{column}{0.5\textwidth}
      \textbf{A pista}\\[0.3em]
      Os cinco contatos aparecem certos no final. Ninguém copiou nada
      explicitamente no seu código.
    \end{column}
  \end{columns}
\end{frame}

\begin{frame}{Definição: \texttt{realloc}}
  \begin{block}{\texttt{void *realloc(void *p, size\_t novos\_bytes);}}
    Redimensiona o bloco apontado por \texttt{p} para \texttt{novos\_bytes} e
    devolve o endereço do bloco resultante --- que pode ser o mesmo ou
    \textbf{outro}. O conteúdo antigo é preservado até o menor dos dois tamanhos.
  \end{block}
  \begin{center}
  \begin{tikzpicture}
    \node[cel] (a0) at (0,0) {ana};
    \node[cel, right=0pt of a0] (a1) {bruno};
    \node[leg, above=2pt of a0.north east] {bloco antigo (2 posições)};

    \node[cel] (b0) at (5.2,0) {ana};
    \node[cel, right=0pt of b0] (b1) {bruno};
    \node[celv, right=0pt of b1] (b2) {?};
    \node[celv, right=0pt of b2] (b3) {?};
    \node[leg, above=2pt of b1.north east] {bloco novo (4 posições)};

    \draw[seta] (a1.east) -- node[above, leg] {copia} node[below, leg] {e libera o antigo} (b0.west);
  \end{tikzpicture}
  \end{center}
  \begin{center}\small
  O espaço acrescentado \textbf{não} é zerado: é indeterminado, como no \texttt{malloc}.
  \end{center}
\end{frame}

\begin{frame}[fragile]{O padrão do ponteiro temporário}
\begin{lstlisting}[basicstyle=\ttfamily\scriptsize]
/* ERRADO: se realloc falhar, devolve NULL e o endereco
   do bloco antigo -- que continua alocado -- se perde. */
v = realloc(v, nova * sizeof(int));

/* CERTO: so' troca v depois de saber que deu certo. */
int *tmp = realloc(v, nova * sizeof(int));
if (tmp == NULL) { /* v ainda vale: da' para se recuperar */ }
v = tmp;
\end{lstlisting}
  \begin{alertblock}{Duas coisas de uma vez}
    A versão errada vaza memória \emph{e} destrói a única cópia do ponteiro para
    os dados que você ainda tinha. Foi por isso que \texttt{agenda.c} usou
    \texttt{n2} e \texttt{i2}.
  \end{alertblock}
\end{frame}

\begin{frame}[fragile]{Armadilha: o ponteiro antigo fica pendurado}
\begin{lstlisting}[basicstyle=\ttfamily\scriptsize]
int *v = malloc(2 * sizeof(int));
v[0] = 10;
int *w = realloc(v, 100 * sizeof(int));
printf("w[0] = %d\n", w[0]);   /* 10: o conteudo veio junto */
printf("v[0] = %d\n", v[0]);   /* v pode ja' ter sido liberado */
\end{lstlisting}
  \begin{alertblock}{O que o AddressSanitizer diz da última linha}
\begin{lstlisting}[basicstyle=\ttfamily\tiny, numbers=none, language=]
ERROR: AddressSanitizer: heap-use-after-free on address 0x6020000000f0
    #0 in main rtrap.c:8
\end{lstlisting}
    Depois de um \texttt{realloc} bem-sucedido, existe \emph{um} ponteiro
    válido: o que ele devolveu.
  \end{alertblock}
\end{frame}

\begin{frame}{Os três papéis de \texttt{realloc}}
  \begin{center}
  \small
  \begin{tabular}{lll}
    \toprule
    Chamada & Efeito & Equivale a \\
    \midrule
    \texttt{realloc(NULL, n)} & aloca do zero              & \texttt{malloc(n)} \\
    \texttt{realloc(p, maior)} & cresce, copiando o conteúdo & --- \\
    \texttt{realloc(p, menor)} & encolhe, descartando o fim  & --- \\
    \bottomrule
  \end{tabular}
  \end{center}
  \begin{exampleblock}{Por que dobrar a capacidade, e não somar 1?}
    Crescer de 1 em 1 até $n$ pode copiar da ordem de $n^2$ elementos.
    Dobrando, o total copiado fica da ordem de $n$: cada elemento é movido, em
    média, um número constante de vezes. Por isso a agenda foi 2, 4, 8.
  \end{exampleblock}
\end{frame}

\begin{frame}{Recapitulando: \texttt{realloc}}
  \begin{block}{O que aprendemos}
    \begin{itemize}
      \item \texttt{realloc} pode devolver um endereço diferente --- e
        frequentemente devolve;
      \item ele copia o conteúdo antigo e libera o bloco antigo por você;
      \item o espaço novo não vem zerado;
      \item receba o resultado em um ponteiro temporário e só então substitua o
        original;
      \item depois da troca, o ponteiro antigo está morto.
    \end{itemize}
  \end{block}
\end{frame}

% ==================================================================
\section{Armadilhas e diagnóstico}
% ==================================================================

\begin{frame}{As quatro falhas clássicas}
  \begin{center}
  \small
  \begin{tabular}{lll}
    \toprule
    Falha & O que é & Sintoma típico \\
    \midrule
    Vazamento          & alocar e nunca liberar        & programa incha com o tempo \\
    Ponteiro pendurado & usar depois do \texttt{free}  & lixo ou \emph{segfault} aleatório \\
    \texttt{free} duplo & liberar duas vezes           & aborta em \texttt{free}, sem pista \\
    Estouro do bloco   & escrever além do tamanho      & corrompe o vizinho \\
    \bottomrule
  \end{tabular}
  \end{center}
  \begin{alertblock}{O que as quatro têm em comum}
    Nenhuma delas falha onde foi cometida. O programa quebra depois, em outro
    lugar, com uma mensagem que aponta para código inocente.
  \end{alertblock}
\end{frame}

\begin{frame}[fragile]{Vazamento: a agenda que esquece os nomes}
\begin{lstlisting}[basicstyle=\ttfamily\scriptsize]
void carrega_agenda(void) {
    char **nomes = malloc(5 * sizeof(char *));
    for (int i = 0; i < 5; i++) {
        nomes[i] = malloc(strlen(entrada[i]) + 1);
        strcpy(nomes[i], entrada[i]);
    }
    free(nomes);   /* liberou o vetor, nao os cinco nomes */
}
\end{lstlisting}
  \begin{alertblock}{Um \texttt{free} não basta}
    \texttt{nomes} guarda cinco endereços. Liberar \texttt{nomes} devolve as
    40 bytes dos \emph{ponteiros} e torna os cinco blocos de texto
    inalcançáveis. Libere de dentro para fora: primeiro cada
    \texttt{nomes[i]}, depois \texttt{nomes}.
  \end{alertblock}
\end{frame}

\begin{frame}[fragile]{Ferramenta 1: o detector de vazamentos}
\begin{lstlisting}[basicstyle=\ttfamily\scriptsize, numbers=none, language=bash]
# macOS
$ MallocStackLogging=1 leaks -atExit -- ./vaza
Process 3103: 5 leaks for 80 total leaked bytes.
STACK OF 5 INSTANCES OF 'ROOT LEAK: <malloc in carrega_agenda>':
      1 (16 bytes) ROOT LEAK: <malloc in carrega_agenda 0xc8100c000>

# Linux
$ gcc -g -fsanitize=address vaza.c -o vaza && ./vaza
$ valgrind --leak-check=full ./vaza
\end{lstlisting}
  \begin{center}\small
  A ferramenta aponta a \emph{função onde o bloco foi alocado} --- e é isso que
  transforma um vazamento em um conserto de duas linhas.
  \end{center}
\end{frame}

\begin{frame}[fragile]{Ferramenta 2: o AddressSanitizer}
\begin{lstlisting}[basicstyle=\ttfamily\scriptsize, numbers=none, language=bash]
$ gcc -Wall -g -fsanitize=address uaf.c -o uaf && ./uaf
ERROR: AddressSanitizer: heap-use-after-free on address 0x6020000000f0
READ of size 2 at 0x6020000000f0 thread T0
    #0 in main uaf.c:8
0x6020000000f0 is located 0 bytes inside of 6-byte region
freed by thread T0 here:
    #0 in free
    #1 in main uaf.c:7
\end{lstlisting}
  \begin{block}{Leia a mensagem inteira}
    Ela diz três coisas: \emph{o que} (uso após liberação), \emph{onde} (linha
    8) e \emph{por causa de quem} (o \texttt{free} da linha 7). Um
    \emph{segfault} comum diria apenas ``o programa morreu''.
  \end{block}
\end{frame}

\begin{frame}{Higiene: cinco hábitos que evitam as quatro falhas}
  \begin{block}{Checklist}
    \begin{enumerate}
      \item pedir \texttt{quantidade * sizeof(tipo)}, nunca um número solto;
      \item testar \texttt{NULL} logo depois de alocar;
      \item receber \texttt{realloc} em um ponteiro temporário;
      \item liberar de dentro para fora, e uma vez só;
      \item \texttt{p = NULL} depois do \texttt{free}, e compilar com
        \texttt{-Wall -g -fsanitize=address} enquanto desenvolve.
    \end{enumerate}
  \end{block}
\end{frame}

\begin{frame}{Recapitulando: armadilhas}
  \begin{block}{O que aprendemos}
    \begin{itemize}
      \item as falhas de memória não quebram o programa onde foram cometidas;
      \item vazar é perder o endereço, não perder os dados --- o bloco continua
        ocupado;
      \item quem alocou em camadas libera em camadas;
      \item as ferramentas transformam ``às vezes dá erro'' em ``linha 8, por
        causa da linha 7''.
    \end{itemize}
  \end{block}
\end{frame}

% ==================================================================
\section{Mãos à obra}
% ==================================================================

\begin{frame}{Altere o programa}
  \begin{block}{Tarefa 1}
    Troque a capacidade inicial de 2 para 1 e o fator de crescimento de
    $\times 2$ para $+1$. Quantas vezes a linha \textbf{C} passa a aparecer?
    Explique o número.
  \end{block}
  \begin{block}{Tarefa 2}
    Troque o \texttt{malloc} de \texttt{idades} por \texttt{calloc} e imprima
    \texttt{idades[7]} logo depois de alocar, antes de preencher. O que muda em
    relação ao \texttt{malloc}?
  \end{block}
  \begin{block}{Tarefa 3}
    Comente o laço que libera \texttt{nomes[i]} e rode o detector de vazamentos.
    Quantos blocos ele acusa? Por que esse número?
  \end{block}
\end{frame}

\begin{frame}{Desafio}
  \begin{block}{Uma agenda que também remove}
    Escreva \texttt{remover(char **nomes, int *idades, int n, const char *quem)}
    que apaga um contato e devolve a nova quantidade.

    \medskip
    Requisitos:
    \begin{itemize}
      \item libere o nome removido antes de perder o endereço dele;
      \item feche o buraco deslocando os contatos seguintes;
      \item se a agenda ficar com menos de um quarto da capacidade, encolha o
        vetor pela metade com \texttt{realloc};
      \item o programa deve terminar sem vazamentos e sem erros do sanitizer.
    \end{itemize}
  \end{block}
\end{frame}

\begin{frame}{Fechamento}
  \begin{columns}[T]
    \begin{column}{0.5\textwidth}
      \textbf{O que mudou hoje}\\[0.3em]
      O tamanho das suas estruturas deixou de ser uma decisão de quem escreve o
      programa e passou a ser uma decisão de quem o executa.
    \end{column}
    \begin{column}{0.5\textwidth}
      \textbf{O preço}\\[0.3em]
      Você ganhou o controle do tempo de vida da memória --- e junto com ele a
      obrigação de devolvê-la, exatamente uma vez, no lugar certo.
    \end{column}
  \end{columns}
  \begin{block}{Para levar}
    \begin{itemize}
      \item \texttt{malloc}: bytes, sem garantia de conteúdo, pode falhar;
      \item \texttt{calloc}: bytes zerados, dois argumentos;
      \item \texttt{realloc}: novo tamanho, endereço possivelmente novo;
      \item \texttt{free}: uma vez, de dentro para fora.
    \end{itemize}
  \end{block}
\end{frame}

\section*{Referências}
\begin{frame}[allowframebreaks]{Referências}
  \nocite{*}
  \bibliographystyle{plain}
  \bibliography{../referencias}
\end{frame}

\end{document}
